\paragraph{Thuật toán Newton--Raphson/IRLS}

Hồi quy Logistic không tồn tại nghiệm dạng đóng cho bài toán tối ưu tham số.
Tuy nhiên, hàm mất mát cross-entropy là hàm lồi và khả vi hai lần, do đó có thể
sử dụng phương pháp Newton--Raphson dựa trên khai triển Taylor bậc hai.

Trong mỗi bước lặp, tham số được cập nhật theo công thức: \begin{equation}
  \boldsymbol{w}^{(t+1)} = \boldsymbol{w}^{(t)} - \boldsymbol{H}^{-1} \nabla E,
\end{equation} trong đó gradient có dạng: \begin{equation}
  \nabla E = \boldsymbol{\Phi}^{\intercal}(\boldsymbol{y} - \boldsymbol{t}),
\end{equation} và Hessian được viết dưới dạng: \begin{equation}
  \boldsymbol{H} = \boldsymbol{\Phi}^{\intercal} \boldsymbol{R}
  \boldsymbol{\Phi} + \lambda \boldsymbol{I},
\end{equation} với: \begin{equation}
  \boldsymbol{R} = \mathrm{diag}\big(y_1(1-y_1), \dots, y_N(1-y_N)\big).
\end{equation}

Thuật toán này còn được gọi là IRLS (Iterative Reweighted Least Squares), do có
thể viết lại dưới dạng bài toán bình phương tối thiểu có trọng số:
\begin{equation}
  \boldsymbol{w}^{(t+1)} =
  (\boldsymbol{\Phi}^{\intercal} \boldsymbol{R} \boldsymbol{\Phi})^{-1}
  \boldsymbol{\Phi}^{\intercal} \boldsymbol{R} \boldsymbol{z},
\end{equation} với: \begin{equation}
  \boldsymbol{z} = \boldsymbol{\Phi}\boldsymbol{w} +
  \boldsymbol{R}^{-1}(\boldsymbol{y}-\boldsymbol{p}).
\end{equation}

IRLS thường hội tụ trong $O(\log(1/\varepsilon))$ bước, nhanh hơn Gradient
Descent cần $O(1/\varepsilon)$ bước. Tuy nhiên, chi phí mỗi bước cao do cần
tính nghịch đảo Hessian, với độ phức tạp xấp xỉ: \begin{equation}
  O(NM^2 + M^3),
\end{equation} trong đó $N$ là số mẫu và $M$ là số đặc trưng.
